% GENERATED FILE. Edit canonical lessons, then run npm run generate:books.
% canonical-source-hash: fnv1a64:44ead4976fa27d6f
% canonical-lessons: UR-C535-shagun, UR-C535-manhus, UR-C535-tona, UR-C535-najumi, UR-C535-zaicha

\chapter{Omens and Stars}
\label{ch:ur-omens-and-stars}
\hlchaptermodality{535}

\begin{chapteropening}
\textbf{By the end of this chapter:} \emph{I can talk about omens and fortune-telling.}

Complete the last lesson of chapter 535: I can talk about omens and fortune-telling.
\end{chapteropening}

\section[\texorpdfstring{\ur{شگون} (shagun)}{شگون (shagun)}]{\ur{شگون} (shagun) — an omen}

\label{lesson:UR-C535-shagun}



\begin{quote}
Before the new one: say the Urdu for checkmate, then the Urdu for ludo.
\end{quote}

\subsection*{You'll want to know: \ur{شگون}}
\textbf{\ur{شگون}} — \emph{shagun} — \textquotedblleft{}an omen\textquotedblright{}.

\begin{quote}
\ur{اگر} \ur{صبح} \ur{بارش} \ur{ہوئی} \ur{تو} \ur{یہ} \ur{اچھا} \ur{شگون} \ur{ہو} \ur{گا۔}
\end{quote}

\begin{quote}
\emph{agar subah bārish huī to yih achchhā shagun ho gā.}
\end{quote}

\begin{quote}
If it rains in the morning, that will be a good omen.
\end{quote}

Masculine.

\begin{cousinweb}
\textbf{\ur{شگون}} (\emph{shagun}) is inherited from Sanskrit \emph{śakuna}, \textquotedblleft{}a bird; an omen\textquotedblright{}, under which CDIAL lists the Hindi cousin \emph{saun}, \textquotedblleft{}a good omen\textquotedblright{} (CDIAL 12245); the fuller form \emph{shagun} is not among its Hindi rows.
\end{cousinweb}

\begin{grammarlens}[title={What you've built}]
One more word for saying what happens if.
\end{grammarlens}

\subsection*{Your turn}
\begin{itemize}
\raggedright
  \item \emph{Say it:} \emph{shagun}
  \item \emph{Say it:} \emph{shagun}, once more
  \item \emph{Say it:} \emph{agar subah bārish huī to yih achchhā shagun ho gā.} — If it rains in the morning, that will be a good omen.
  \item \emph{Recall:} say the Urdu for checkmate, then the Urdu for ludo, then say \emph{shagun} again
\end{itemize}

\begin{tcolorbox}[breakable,colback=teal!4,colframe=teal!35!black,title={Before you move on}]
What does \ur{شگون} mean? (\textquotedblleft{}An omen\textquotedblright{}.) Say it once more.
\end{tcolorbox}
\section[\texorpdfstring{\ur{منحوس} (manhūs)}{منحوس (manhūs)}]{\ur{منحوس} (manhūs) — unlucky; ill-omened}

\label{lesson:UR-C535-manhus}



\begin{quote}
Before the new one: say the Urdu for ludo, then the Urdu for an omen.
\end{quote}

\subsection*{You'll want to know: \ur{منحوس}}
\textbf{\ur{منحوس}} — \emph{manhūs} — \textquotedblleft{}unlucky; ill-omened\textquotedblright{}.

\begin{quote}
\ur{اگر} \ur{یہ} \ur{گھر} \ur{منحوس} \ur{ہوتا} \ur{تو} \ur{ہم} \ur{یہاں} \ur{نہ} \ur{رہتے۔}
\end{quote}

\begin{quote}
\emph{agar yih ghar manhūs hotā to ham yahāṅ na rahte.}
\end{quote}

\begin{quote}
If this house were unlucky, we wouldn't live here.
\end{quote}

\begin{cousinweb}
\textbf{\ur{منحوس}} (\emph{manhūs}) is usually traced, through Persian, to Arabic \emph{manḥūs}, \textquotedblleft{}ill-starred, unlucky\textquotedblright{}.
\end{cousinweb}

\begin{grammarlens}[title={What you've built}]
One more word for saying what happens if.
\end{grammarlens}

\subsection*{Your turn}
\begin{itemize}
\raggedright
  \item \emph{Say it:} \emph{manhūs}
  \item \emph{Say it:} \emph{manhūs}, once more
  \item \emph{Say it:} \emph{agar yih ghar manhūs hotā to ham yahāṅ na rahte.} — If this house were unlucky, we wouldn't live here.
  \item \emph{Recall:} say the Urdu for ludo, then the Urdu for an omen, then say \emph{manhūs} again
\end{itemize}

\begin{tcolorbox}[breakable,colback=teal!4,colframe=teal!35!black,title={Before you move on}]
What does \ur{منحوس} mean? (\textquotedblleft{}Unlucky; ill-omened\textquotedblright{}.) Say it once more.
\end{tcolorbox}
\section[\texorpdfstring{\ur{ٹونا} (ṭonā)}{ٹونا (ṭonā)}]{\ur{ٹونا} (ṭonā) — a magic spell}

\label{lesson:UR-C535-tona}



\begin{quote}
Before the new one: say the Urdu for an omen, then the Urdu for ill-omened.
\end{quote}

\subsection*{You'll want to know: \ur{ٹونا}}
\textbf{\ur{ٹونا}} — \emph{ṭonā} — \textquotedblleft{}a magic spell\textquotedblright{}.

\begin{quote}
\ur{اگر} \ur{ٹونا} \ur{سچ} \ur{ہوتا} \ur{تو} \ur{سب} \ur{امیر} \ur{ہوتے۔}
\end{quote}

\begin{quote}
\emph{agar ṭonā sach hotā to sab amīr hote.}
\end{quote}

\begin{quote}
If spells were real, everyone would be rich.
\end{quote}

Masculine.

\begin{cousinweb}
\textbf{\ur{ٹونا}} (\emph{ṭonā}) is inherited: CDIAL lists the Hindi cousin \emph{ṭonā}, \textquotedblleft{}sorcery\textquotedblright{}, under \emph{*ṭōna}, a form it reconstructs and perhaps links with \emph{śakuna}, \textquotedblleft{}a bird of omen\textquotedblright{} (CDIAL 5480).
\end{cousinweb}

\begin{grammarlens}[title={What you've built}]
One more word for saying what happens if.
\end{grammarlens}

\subsection*{Your turn}
\begin{itemize}
\raggedright
  \item \emph{Say it:} \emph{ṭonā}
  \item \emph{Say it:} \emph{ṭonā}, once more
  \item \emph{Say it:} \emph{agar ṭonā sach hotā to sab amīr hote.} — If spells were real, everyone would be rich.
  \item \emph{Recall:} say the Urdu for an omen, then the Urdu for ill-omened, then say \emph{ṭonā} again
\end{itemize}

\begin{tcolorbox}[breakable,colback=teal!4,colframe=teal!35!black,title={Before you move on}]
What does \ur{ٹونا} mean? (\textquotedblleft{}A magic spell\textquotedblright{}.) Say it once more.
\end{tcolorbox}
\section[\texorpdfstring{\ur{نجومی} (najūmī)}{نجومی (najūmī)}]{\ur{نجومی} (najūmī) — an astrologer}

\label{lesson:UR-C535-najumi}



\begin{quote}
Before the new one: say the Urdu for ill-omened, then the Urdu for a magic spell.
\end{quote}

\subsection*{You'll want to know: \ur{نجومی}}
\textbf{\ur{نجومی}} — \emph{najūmī} — \textquotedblleft{}an astrologer\textquotedblright{}.

\begin{quote}
\ur{اگر} \ur{نجومی} \ur{سب} \ur{جانتا} \ur{تو} \ur{وہ} \ur{خود} \ur{امیر} \ur{ہوتا۔}
\end{quote}

\begin{quote}
\emph{agar najūmī sab jāntā to voh khud amīr hotā.}
\end{quote}

\begin{quote}
If the astrologer knew everything, he would be rich himself.
\end{quote}

Masculine or feminine, by the person meant.

\begin{cousinweb}
\textbf{\ur{نجومی}} (\emph{najūmī}) is usually traced, through Persian, to Arabic \emph{nujūm}, \textquotedblleft{}stars\textquotedblright{}, the plural of \emph{najm}, with the ending \emph{-ī}, \textquotedblleft{}one who\textquotedblright{}.
\end{cousinweb}

\begin{grammarlens}[title={What you've built}]
One more word for saying what happens if.
\end{grammarlens}

\subsection*{Your turn}
\begin{itemize}
\raggedright
  \item \emph{Say it:} \emph{najūmī}
  \item \emph{Say it:} \emph{najūmī}, once more
  \item \emph{Say it:} \emph{agar najūmī sab jāntā to voh khud amīr hotā.} — If the astrologer knew everything, he would be rich himself.
  \item \emph{Recall:} say the Urdu for ill-omened, then the Urdu for a magic spell, then say \emph{najūmī} again
\end{itemize}

\begin{tcolorbox}[breakable,colback=teal!4,colframe=teal!35!black,title={Before you move on}]
What does \ur{نجومی} mean? (\textquotedblleft{}An astrologer\textquotedblright{}.) Say it once more.
\end{tcolorbox}
\section[\texorpdfstring{\ur{زائچہ} (zāichā)}{زائچہ (zāichā)}]{\ur{زائچہ} (zāichā) — a horoscope}

\label{lesson:UR-C535-zaicha}



\begin{quote}
Before the new one: say the Urdu for a magic spell, then the Urdu for an astrologer.
\end{quote}

\subsection*{You'll want to know: \ur{زائچہ}}
\textbf{\ur{زائچہ}} — \emph{zāichā} — \textquotedblleft{}a horoscope\textquotedblright{}.

\begin{quote}
\ur{اگر} \ur{نجومی} \ur{نے} \ur{زائچہ} \ur{بنایا} \ur{تو} \ur{وہ} \ur{پیسے} \ur{مانگے} \ur{گا۔}
\end{quote}

\begin{quote}
\emph{agar najūmī ne zāichā banāyā to voh paise māṅge gā.}
\end{quote}

\begin{quote}
If the astrologer draws up a horoscope, he will ask for money.
\end{quote}

Masculine.

\begin{cousinweb}
\textbf{\ur{زائچہ}} (\emph{zāichā}) is usually traced to Persian \emph{zāyicha}, \textquotedblleft{}a horoscope\textquotedblright{}, a chart of the stars at a birth.
\end{cousinweb}

\begin{grammarlens}[title={What you've built}]
One more word for saying what happens if.
\end{grammarlens}

\subsection*{Your turn}
\begin{itemize}
\raggedright
  \item \emph{Say it:} \emph{zāichā}
  \item \emph{Say it:} \emph{zāichā}, once more
  \item \emph{Say it:} \emph{agar najūmī ne zāichā banāyā to voh paise māṅge gā.} — If the astrologer draws up a horoscope, he will ask for money.
  \item \emph{Recall:} say the Urdu for a magic spell, then the Urdu for an astrologer, then say \emph{zāichā} again
\end{itemize}

\begin{tcolorbox}[breakable,colback=teal!4,colframe=teal!35!black,title={Before you move on}]
What does \ur{زائچہ} mean? (\textquotedblleft{}A horoscope\textquotedblright{}.) Say it once more.
\end{tcolorbox}
